\section{Results}
Greenhalgh et al. reported that structural-variation burden varied by
roughly two orders of magnitude across pediatric cancer types,
underscoring that rearrangement biology is highly disease-specific
rather than uniform across childhood malignancies \supercite{Greenhalgh_2026}.
Compared with adult brain and solid tumors, most pediatric solid cancers
carried substantially fewer structural variants, whereas hematological
malignancies showed a much closer burden profile.

To make formula support easier to review by eye, the demo manuscript now
includes a compact set of representative expressions. Let $\rho$ denote
a rank-based association summary, let $\Delta$ denote a between-group
burden shift, let $\hat{\lambda}$ denote an estimated scaling factor,
and let $n_{\mathrm{SV},i}$ denote the structural-variation count for
sample $i$. A log-scaled burden summary can then be written as
\[
B_i = \log_{10}\!\left(n_{\mathrm{SV},i} + 1\right).
\]
Using cohort-level mean $\mu_B$ and standard deviation $\sigma_B$, the
standardized burden score is
\[
z_i = \frac{B_i - \mu_B}{\sigma_B}.
\]
If multiple signature-derived components are combined into a single
review score, one convenient weighted summary is
\[
S_i = \sum_{k=1}^{K} w_k x_{ik},
\qquad
\sum_{k=1}^{K} w_k = 1.
\]
For a simple threshold-based flag, the indicator variable can be
expressed as
\[
I_i =
\begin{cases}
1, & z_i \geq 2, \\
0, & \text{otherwise}.
\end{cases}
\]
Together, these examples cover inline Greek letters, subscripts,
superscripts, fractions, summation, delimiters, and a piecewise
definition in both PDF and DOCX export.

The same analysis ranked genes disrupted by structural variants and
found that the highest-priority events aligned with recognized pediatric
cancer drivers, while adult cohorts were enriched for disruptions
involving fragile genomic regions \supercite{Greenhalgh_2026}. This contrast suggested
that pediatric structural variation more often tracks direct driver
selection, whereas adult rearrangement landscapes carry a larger imprint
of accumulated genomic instability.

A second notable result concerned hotspot architecture near RAG
recombination signal sequences. In pediatric acute lymphoblastic
leukemias, recurrent hotspots affected both immune loci and oncogenic
driver genes, while adult lymphoid cancers were described as showing
disruption of immune loci without the same degree of pediatric driver
involvement \supercite{Greenhalgh_2026}.

Finally, signature extraction and longitudinal clustering added an
evolutionary perspective. The authors described ten structural-variation
signatures, implicated RAG-mediated mutagenesis in the etiology of
COSMIC SV7 within lymphoid cancers, and used multiple samples from 13
patients to illustrate that structural variants continue to shape
intra-tumor heterogeneity over time \supercite{Greenhalgh_2026}.
